\documentclass[12pt, a4paper]{article}

% --- Pacotes Básicos ---
\usepackage[utf8]{inputenc}
\usepackage[T1]{fontenc}
\usepackage[main=brazilian, english]{babel}
\usepackage{indentfirst}
\usepackage{geometry}
\geometry{left=3cm, top=3cm, right=2cm, bottom=2cm}

% --- Informações do Documento ---
\title{Resenha: O Preço da Verdade}
\author{Erickson Giesel Müller}
\date{\today}

\begin{document}
\maketitle
\thispagestyle{empty} % primeira página não tem número.

	O filme "O preço da verdade" trata a respeito da história de um advogado corporativo que defende as causas de uma empresa química quando um fazendeiro, amigo de sua avó, entra em contato buscando apoio jurídico para processar a corporação que está contaminando a água da cidade, resultando na morte de seus animais. Esse caso coloca em jogo a moral do advogado, que construiu a sua carreira defendendo as corporações químicas e agora precisa colocar ignorar seu "instinto de advogado" de defender seu cliente não importa seu crime e usar esse conhecimento no novo caso.

	A troca de perspectiva levou o protagonista à loucura, ele passa boa parte do filme buscando informações sobre o caso de uma maneira maníaca que o coloca, diversas vezes, em posição de humilhação perante seus colegas. Em determinado momento, o advogado entra em contato com um químico para que este explique os efeitos colaterais da contaminação dos produtos químicos que a empresa está soltando na água. O personagem percebe que as ações da empresa põe em risco não apenas os animais da fazenda de seu cliente, mas a vida de toda a população da cidade, de sua avó e do seu filho que está para nascer.

	Sob a ótica da disciplina, a narrativa escancara uma grave violação de garantias fundamentais, em especial o direito ao meio ambiente e o direito à saúde. A postura da corporação ao longo de décadas reflete a priorização irrestrita do lucro em detrimento da dignidade humana, transformando os moradores da cidade em vítimas invisíveis que não tem força para se defenderem. O filme ilustra como os interesses econômicos hegemônicos podem se sobrepor à responsabilidade social, atropelando a base de direitos que deveriam ser inegociáveis para qualquer cidadão. Percebemos isso na cena em que o fazendeiro e o advogado discutem sobre a impossibilidade de trazer um cientista para o caso, pois qualquer químico que entenda a respeito do material contaminante em questão já estaria contratado por essas empresas químicas.

Além disso, a história levanta um debate crucial sobre a omissão estatal no que diz respeito à proteção da coletividade. A falha das agências reguladoras em fiscalizar e frear os abusos da indústria química evidencia um grave ato de corrupção, onde o interesse público acaba sendo atropelado pela força do lobby corporativo. Essa inércia do Estado não apenas deixa a população desamparada, mas obriga a sociedade civil a buscar seus direitos básicos por meios burocráticos e inacessíveis.

Por fim, "O Preço da Verdade" é uma mensagem poderosa sobre a importância do acesso à justiça e do exercício ativo da cidadania. A exaustiva jornada do protagonista, aliada à resistência e ao sofrimento da comunidade local, demonstra que a cidadania não é uma condição estática, mas uma prática contínua de reivindicação e enfrentamento, vemos que a posição natural da população é, infelizmente, sob a opressão violenta de pessoas mais poderosas.


\end{document}
